% Complete chapter-2 typesetting source; compact textbook presentation.
% Authoritative Markdown: ../../chapters/ch02/student.md
% Stable task/source IDs are comments. All prose, tasks, tables, and notes retained.
% Local editorial-file links retain their labels; external links remain live.
% Shared layout for the two complete chapter-2 PDFs.
% This is a small hand-maintained TeX layout, not a conversion pipeline.
\RequirePackage{fix-cm}
\documentclass[11pt,twoside]{article}
\usepackage{fontspec}
\setmainfont{Cambria}[
  Renderer=HarfBuzz,
  BoldFont=cambriab.ttf,
  ItalicFont=cambriai.ttf,
  BoldItalicFont=cambriaz.ttf
]
\setmonofont{lmmono10-regular.otf}[BoldFont=lmmonolt10-bold.otf]
\usepackage[paperwidth=165mm,paperheight=241mm,inner=21mm,outer=19mm,
  top=22mm,bottom=20mm,headheight=14pt,headsep=7mm,footskip=10mm]{geometry}
\usepackage{xcolor,array,booktabs,longtable,tabularx}
\usepackage{enumitem,fancyhdr}
\usepackage[unicode,hidelinks]{hyperref}
\definecolor{ink}{HTML}{000000}
\definecolor{muted}{HTML}{4C5558}
\setlength{\parindent}{0pt}
\setlength{\parskip}{3.5pt}
\AtBeginDocument{\fontsize{11}{14.3}\selectfont}
\setlength{\tabcolsep}{4pt}
\renewcommand{\arraystretch}{1.10}
\setlist[itemize]{leftmargin=1.3em,label=\textbullet,itemsep=2pt,topsep=3pt,parsep=0pt}
\setlist[enumerate]{leftmargin=1.8em,itemsep=3pt,topsep=4pt,parsep=0pt}
\widowpenalty=10000
\clubpenalty=10000
\emergencystretch=1em
\raggedbottom
\pagestyle{fancy}
\fancyhf{}
\fancyhead[LE,RO]{\small\itshape Plural nouns}
\fancyhead[LO,RE]{\small Cook Expanded}
\fancyfoot[LE,RO]{\small\thepage}
\renewcommand{\headrulewidth}{0.3pt}
\newcommand{\booktitle}[2]{
  {\fontsize{9}{11}\selectfont\color{muted}#1\par}
  \vspace{3pt}
  {\fontsize{20}{24}\selectfont\color{ink}#2\par}
  \vspace{6pt}
}
\newcommand{\source}[1]{\par{\footnotesize\color{muted}#1\par}}
\newcommand{\readingtext}[1]{
  \par\addvspace{5pt}
  \begingroup\leftskip=1.5em
    \fontsize{12}{16}\selectfont #1\par
  \endgroup\addvspace{5pt}
}
\makeatletter
\renewcommand{\section}{\@startsection{section}{1}{\z@}
  {-2.4ex \@plus -0.6ex \@minus -0.2ex}{0.8ex \@plus 0.2ex}
  {\normalfont\fontsize{12}{15}\selectfont\bfseries\color{ink}}}
\renewcommand{\subsection}{\@startsection{subsection}{2}{\z@}
  {-1.8ex \@plus -0.5ex \@minus -0.2ex}{0.5ex \@plus 0.2ex}
  {\normalfont\fontsize{11}{14}\selectfont\bfseries\color{ink}}}
\makeatother

\hypersetup{pdftitle={Chapter 2: Plural nouns},pdfauthor={Cook Expanded}}
\fancyfoot[LO,RE]{\footnotesize\color{muted}Chapter 2}
\begin{document}

\booktitle{CHAPTER 2}{Plural nouns}

% Chapter: ch02. Independently reviewed. Source crosswalk: coverage.json.

\pdfbookmark[1]{1. Nominative plural}{chapter-section-1}
\section*{1. Nominative plural}


The nominative plural article is \textbf{ðā} for all three genders. A noun retains its gender in the plural.


% example: ch02-example-01

\textbf{se fisc \ensuremath{\rightarrow} ðā fiscas}; \textbf{sēo wyrt \ensuremath{\rightarrow} ðā wyrta}; \textbf{ðæt dēor \ensuremath{\rightarrow} ðā dēor}.


Plural formation depends on declension, not gender alone. The following patterns occur in this lesson:


\begin{longtable}{@{}>{\raggedright\arraybackslash}p{55mm}>{\raggedright\arraybackslash}p{67mm}@{}}
\toprule
\textbf{Pattern} & \textbf{Example} \\
\midrule
\endfirsthead
\toprule
\textbf{Pattern} & \textbf{Example} \\
\midrule
\endhead
\midrule
\multicolumn{2}{r}{\footnotesize\itshape Continued on the next page} \\
\endfoot
\bottomrule
\endlastfoot
Masculine \textbf{-as} & \textbf{fisc \ensuremath{\rightarrow} fiscas} \\
Feminine \textbf{-a} & \textbf{wyrt \ensuremath{\rightarrow} wyrta} \\
Neuter \textbf{-u} & \textbf{nīeten \ensuremath{\rightarrow} nīetenu} \\
Neuter, no added ending & \textbf{dēor \ensuremath{\rightarrow} dēor} \\
\end{longtable}


Stem changes accompany some endings: \textbf{dæg \ensuremath{\rightarrow} dagas}, \textbf{hwæl \ensuremath{\rightarrow} hwalas}; \textbf{fugol \ensuremath{\rightarrow} fuglas}, \textbf{morgen \ensuremath{\rightarrow} morgnas}, \textbf{tācen \ensuremath{\rightarrow} tācnu}. \textbf{trēo} has \textbf{w} before the ending: \textbf{trēowu}. Nouns in \textbf{-nes} have \textbf{-nessa}. Learn \textbf{niht}, \textbf{męnn}, and \textbf{sǣs} individually.


\textbf{Neuter -u and zero.} The neuters in §2 belong to the strong declension. In monosyllabic strong neuters, a light stem—a short vowel or short diphthong followed by one consonant—normally takes \textbf{-u}. A heavy stem—a long vowel or long diphthong, or a short vowel or diphthong followed by two or more consonants (including a doubled consonant)—normally has no ending. Thus \textbf{gēar, dēor, sǣd} remain unchanged.

% example: ch02-example-02

Supplied illustrations, not vocabulary for recall: \textbf{hof \ensuremath{\rightarrow} hofu}, dwelling(s); \textbf{word \ensuremath{\rightarrow} word}, word(s).

Two-syllable nouns vary; a heavy first syllable followed by a light second syllable often gives \textbf{-u}: \textbf{nīetenu, ǣfenu}. Some lose the second vowel before \textbf{-u}: \textbf{tācen \ensuremath{\rightarrow} tācnu}. Retain the forms in §2. \textbf{trēo \ensuremath{\rightarrow} trēowu} follows a separate pattern despite its long diphthong.

The nominative plural is used for plural subjects and for the isolated phrases below. Coordinate definite phrases as before: article + noun, \textbf{and}, article + noun.


\clearpage
\pdfbookmark[1]{2. Vocabulary}{chapter-section-2}
\section*{2. Vocabulary}


All forms are nominative. New headwords: \textbf{sǣ}, \textbf{hwæl}, \textbf{nīeten}. The other headwords are familiar; all twenty-two plural forms are new.


\begin{longtable}{@{}>{\raggedright\arraybackslash}p{25mm}>{\raggedright\arraybackslash}p{16mm}>{\raggedright\arraybackslash}p{29mm}>{\raggedright\arraybackslash}p{46.5mm}@{}}
\toprule
\textbf{Singular} & \textbf{Gender} & \textbf{Plural} & \textbf{Meaning} \\
\midrule
\endfirsthead
\toprule
\textbf{Singular} & \textbf{Gender} & \textbf{Plural} & \textbf{Meaning} \\
\midrule
\endhead
\midrule
\multicolumn{4}{r}{\footnotesize\itshape Continued on the next page} \\
\endfoot
\bottomrule
\endlastfoot
\textbf{dæg} & m. & \textbf{dagas} & day \\
\textbf{niht} & f. & \textbf{niht} & night \\
\textbf{fisc} & m. & \textbf{fiscas} & fish \\
\textbf{fugol} & m. & \textbf{fuglas} & bird \\
\textbf{trēo} & n. & \textbf{trēowu} & tree \\
\textbf{wyrt} & f. & \textbf{wyrta} & herb; plant \\
\textbf{andlīcnes} & f. & \textbf{andlīcnessa} & image \\
\textbf{gelīcnes} & f. & \textbf{gelīcnessa} & likeness \\
\textbf{tācen} & n. & \textbf{tācnu} & sign \\
\textbf{tīd} & f. & \textbf{tīda} & season; time \\
\textbf{gēar} & n. & \textbf{gēar} & year \\
\textbf{mann} & m. & \textbf{męnn} & man; human being \\
\textbf{dēor} & n. & \textbf{dēor} & animal; beast \\
\textbf{sǣ} & m. & \textbf{sǣs} & sea \\
\textbf{heofon} & m. & \textbf{heofonas} & heaven; sky \\
\textbf{sǣd} & n. & \textbf{sǣd} & seed \\
\textbf{wæstm} & m. & \textbf{wæstmas} & fruit; produce \\
\textbf{ǣfen} & n. & \textbf{ǣfenu} & evening \\
\textbf{morgen} & m. & \textbf{morgnas} & morning \\
\textbf{gāst} & m. & \textbf{gāstas} & spirit \\
\textbf{hwæl} & m. & \textbf{hwalas} & whale \\
\textbf{nīeten} & n. & \textbf{nīetenu} & animal; beast; in the plural, cattle \\
\end{longtable}


\pdfbookmark[1]{3. Preparation}{chapter-section-3}
\section*{3. Preparation}


\begin{enumerate}
% item: ch02-prep-01
\item[1.] Put into the nominative plural: \textbf{se dæg}; \textbf{sēo wyrt}; \textbf{ðæt dēor}.
% item: ch02-prep-02
\item[2.] Give the nominative plural of \textbf{fugol}, \textbf{tācen}, \textbf{mann}.
% item: ch02-prep-03
\item[3.] Give the meaning and gender of \textbf{sǣ}, \textbf{hwæl}, \textbf{nīeten}.
\end{enumerate}


\clearpage
\pdfbookmark[1]{4. Composition}{chapter-section-4}
\section*{4. Composition}


Translate into Old English. Use the vocabulary of §2, nominative plural forms, an article before each noun, and \textbf{and}.


\begin{enumerate}
% item: EX-I-02-01
\item[1.] The days and the nights.
% item: EX-I-02-02
\item[2.] The fishes and the birds.
% item: EX-I-02-03
\item[3.] The trees and the herbs.
% item: EX-I-02-04
\item[4.] The images and the likenesses.
% item: EX-I-02-05
\item[5.] The signs and the seasons.
% item: EX-I-02-06
\item[6.] The years and the days.
% item: EX-I-02-07
\item[7.] The men and the animals.
% item: EX-I-02-08
\item[8.] The seas and the heavens.
% item: EX-I-02-09
\item[9.] The seeds and the fruits.
% item: EX-I-02-10
\item[10.] The evenings and the mornings.
% item: EX-I-02-11
\item[11.] The spirits and the men.
% item: EX-I-02-12
\item[12.] The whales and the cattle.
\end{enumerate}


\clearpage
\pdfbookmark[1]{5. Reading and translation}{chapter-section-5}
\section*{5. Reading and translation}


A clause from God's blessing of the living creatures in the Creation narrative.


The following forms are supplied, not assigned for independent recall.


\begin{longtable}{@{}>{\raggedright\arraybackslash}p{34mm}>{\raggedright\arraybackslash}p{88mm}@{}}
\toprule
\textbf{Form} & \textbf{Meaning and analysis} \\
\midrule
\endfirsthead
\toprule
\textbf{Form} & \textbf{Meaning and analysis} \\
\midrule
\endhead
\midrule
\multicolumn{2}{r}{\footnotesize\itshape Continued on the next page} \\
\endfoot
\bottomrule
\endlastfoot
\textbf{bēon} & be; present plural optative (subjunctive), here expressing a command or wish: “let … be.” \\
\textbf{gemanigfielde} & multiplied; increased in number. Past participle, masculine nominative plural. \\
\textbf{ofer} & over; across. Preposition. \\
\textbf{eorðan} & earth; oblique singular of \textbf{eorðe} after \textbf{ofer}. \\
\end{longtable}


\textbf{bēon gemanigfielde}, “let … be multiplied”; \textbf{ofer eorðan}, “over/across the earth.”


% item: ch02-text-01

\readingtext{and ðā fuglas bēon gemanigfielde ofer eorðan.}


Using the supplied forms, answer in English:


\begin{enumerate}
% item: ch02-reading-01
\item[1.] What is to increase, and where?
% item: ch02-reading-02
\item[2.] Identify the subject's article and noun; give gender, number, and case.
% item: ch02-reading-03
\item[3.] Translate the clause.
\end{enumerate}


\pdfbookmark[1]{6. Review}{chapter-section-6}
\section*{6. Review}


For 1–2, use the taught vocabulary, nominative plural forms, an article before each noun, and \textbf{and}.


\begin{enumerate}
% item: ch02-review-01
\item[1.] The seeds and the birds.
% item: ch02-review-02
\item[2.] The whales and the seas.
% item: ch02-review-03
\item[3.] Give the number of each phrase and explain the distinction: \textbf{ðæt sǣd}; \textbf{ðā sǣd}. Why does \textbf{sǣd} have no plural ending?
\end{enumerate}


\clearpage
\begingroup\small
\pdfbookmark[1]{Notes on sources and editorial choices}{chapter-section-7}
\section*{Notes on sources and editorial choices}


\textbf{Ex:} A. S. Cook, \emph{Exercises in Old English} (1895). \textbf{FB:} \emph{A First Book in Old English}, 2nd ed. (1895). References use printed pages; full evidence and scan locators are in the editorial record.


\begin{enumerate}
\item[1.] \textbf{Grammar and vocabulary (§§1–2):} FB §§4, 17, 43, 46–47, 51–52, 84. Stem weight: §4, p.6; neuter endings and the supplied \textbf{hof/word} paradigms: §47, pp.32–33; disyllable variation, \href{https://doi.org/10.1017/S1360674321000435}{Rodríguez Ledesma (2022), §5}. New headwords: Ex, s.v. cattle, p.42; sea, p.60; whale, p.67; \textbf{nīeten}, FB vocabulary p.307. \textbf{tācnu} and \textbf{morgnas} are corroborated by \href{https://bosworthtoller.com/58972}{Bosworth–Toller, \emph{tácn}} and \href{https://bosworthtoller.com/23129}{\emph{morgen}}. \textbf{ǣfenu} follows the neuter patterns in FB §47.1, 7; it is a grammatical formation, not a quoted occurrence. The key records alternatives to the selected forms.
\item[2.] \textbf{Correction C01 (§2; composition 9):} Short \textbf{æ} in \textbf{wæstm, wæstmas} continues Chapter 1's correction of Cook's \textbf{wǣstm}: \href{https://glossary.oldenglishaerobics.net/oeaginfo.php?objid=xwaestm\&objtype=gentry}{Baker, \emph{wæstm}}; \href{https://doi.org/10.17811/selim.29.2024.109-25}{Stenroos (2024), p.118, Table 1}.
\item[3.] \textbf{Exercises (§§3–6):} All twelve composition prompts preserve Ex, Group I, Exercise 2, p.1. Directions, examples, additional tasks, and model answers are new teaching material.
\item[4.] \textbf{Reading (§5):} FB, Reader I, p.126, printed lines 12–13. The extracted clause retains every word and vowel mark; cue 21 and the closing quotation mark belonging to the surrounding speech are omitted. \textbf{bēon}: §138, p.78; §193(a), p.103. \textbf{gemanigfielde}: p.126, note 19, and vocabulary p.293. The translation is independently prepared. The case of \textbf{eorðan} is not tested; the form permits accusative and dative analyses after \textbf{ofer}.
\end{enumerate}

\endgroup

\end{document}
